\begin{table}[htbp]
\centering
\footnotesize
\caption{各设定下相对损失的截距 $a$、斜率 $b$ 与平衡比 $b/a$，以及各估计值的双向自助法 95\% 置信区间。}
\label{tab:s1}
\begin{tabular}{lrrrlllr}
\toprule
设定 & $N$ & 周数 & $\bar{x}$ & $a$ [95\% CI] & $b$ [95\% CI] & $b/a$ [95\% CI] & 标准差比 \\
\midrule
$R_A$，规则窗口（$S=4$） & 1748 & 42 & 1.206 & 0.0150 [0.0100, 0.0208] & 0.560 [0.367, 0.769] & 37.4 [19.1, 71.0] & 1.95 \\
$G$，规则窗口（$S=4$） & 1756 & 42 & 1.724 & 0.0840 [0.0498, 0.1248] & 1.971 [0.959, 3.098] & 23.5 [8.4, 59.7] & 2.29 \\
$G$，日历窗口（$S=3$） & 1267 & 29 & 1.671 & 0.0991 [0.0569, 0.1382] & 2.243 [1.060, 3.559] & 22.6 [9.0, 52.2] & 4.68 \\
$R_A$，日历窗口（$S=3$） & 1263 & 29 & 1.190 & 0.0202 [0.0097, 0.0306] & 0.516 [0.279, 0.839] & 25.6 [10.2, 79.6] & 3.46 \\
\bottomrule
\end{tabular}
\par\smallskip
\begin{minipage}{0.95\linewidth}\scriptsize 以迭代 FGLS 将 $L=((x-\bar{x})/\bar{x})^2$ 拟合为 $a+b/D$；区间为州 $\times$ 周双向鸽笼自助法 95\% 百分位区间，$B=2000$，每次重抽重新估计 $\bar{x}$。标准差比为双向重抽与仅按州重抽下 $a$ 的自助标准差之比。均为 51 个辖区。\end{minipage}
\end{table}
